% Page 027 — page imprimée 23
% Contenu seul : le préambule et la pagination sont gérés par meyrueis.tex

\textbf{Il avait donc fallu 6 mois pour qu’un membre du Comité Populaire se soucie du sort de
la femme du pasteur Samuel François, et ose, suivi par tous les membres du Comité
Populaire de Meyrueis, se dresser contre la mesure unique prise par Châteauneuf-
Randon. !!!}

Quelques mois plus tard, les poursuites contre Samuel François sont levées. Il peut reparaître
en plein jour, et regagne Meyrueis, sa famille et les séances de la Société Populaire. Voilà
comment est relaté son retour (séance du 21 ventôse An III (11 mars 1795).

\begin{quote}
\itshape
«\ldots\ Un sujet intéressant devait fixer dans cette séance l’attention et l’attendrissement de
chaque citoyen : c’était le récit de la persécution affreuse à laquelle a été en butte toute une
année, l’honnête citoyen dont tout le crime était d’avoir osé défendre les droits du peuple. Il a
demandé d’être réintégré dans la Société ; il a même produit des preuves de son civisme
reconnu pendant sa fuite ; mais l’assemblée assez pénétrée depuis longtemps des vertus
énergiques du citoyen Samuel François, l’a admis comme par l’acclamation et l’accolade
fraternelle qu’il a reçue des membres du bureau, lui est un sûr garant de la joie que ressentait
la Société du recouvrement d’un membre, dont l’absence n’avait pu, que se faire sentir
vivement dans son sein et dans l’administration du Département, dont il était membre. »
\end{quote}

\textbf{En mars 1794 personne n’avait osé, au moins en public, prendre la défense de Samuel
François alors qu’il venait de s’évader : ce n’était qu’un scélérat, mais un an plus tard,
en Mars 1795, les membres de la Société se jettent à son cou comme si de rien n’était.}

Il est vrai que \textbf{Châteauneuf-Randon} (\textit{de son vrai nom Alexandre Paul Guérin du
Tournel, marquis de Joyeuse, Comte de Châteauneuf-Randon}) et ses méthodes violentes
ne sévissait plus en Lozère. (Il lui fallait bien faire la preuve, lui, l’aristocrate, qu’il était un
vrai montagnard -- révolutionnaire). Il avait été \textbf{nommé auprès de l’armée des Pyrénées}.

\vfill

\begin{center}
\small
Extraits des comptes-rendus de la Société Populaire de Meyrueis,\\
édités par la Société d’Agriculture, Industrie, Sciences et\\
Arts de la Lozère, 1913. Archives Gévaudanaises\\
Mars 2006
\end{center}
